\section{The Economic Decision and Its Robustness}
\label{sec:decision_target}
The object of a portfolio certificate is a decision under a specified contract, not a sign attached to an unspecified limiting model. Three primitives are particularly important here: the duration for which the initial position is chosen, the settlement technology that determines subsequent states, and the cardinal value assigned to a change in preferences. We make these primitives explicit and study their effects without changing the financial controls, adjustment technology, or liquidation rule of the baseline experiment.

\subsection{A Rebalancing Contract with a Finite Settlement Lottery}
\label{subsec:finite_protocol}
The finite economy has eight rebalancing dates over one year. At each date the agent observes $(u,X)$, chooses $(c,\theta,\pi)$, and holds that control until the next date or liquidation, whichever occurs first. The preference grid contains 33 tiers on $[1.2,2.8]$; the baseline wealth grid contains 49 balances on $[0.5,2]$. The union of the two action meshes has 1,565 elements. Three deposited state- and date-dependent feasible proposals complete the 1,568-action baseline menu. Removing those proposals from a lower-policy construction does not remove them from the target against which that policy is certified.

Here is a primitive description of the transition experiment. Independently at every date, a law indicator chooses the correlation $-0.25$ with probability $1-\lambda$ and $0.25$ with probability $\lambda$. The probability $\lambda$ is fixed for the entire contract; the indicator is not observed before the action. Conditional on the indicator, two independent symmetric signs determine a proposed increment using the drift and volatility in \eqref{eq:ndu_u}--\eqref{eq:ndu_x}. For an interval $h=1/8$ and signs $(z_u,z_X)$, that increment is
\begin{equation}
 h\begin{pmatrix}\theta\\(r+\pi(\mu-r))X-c\end{pmatrix}
 +\sqrt h\begin{pmatrix}\sigma_u z_u\\
 \pi X\sigma_X\{\varrho z_u+\sqrt{1-\varrho^2}z_X\}\end{pmatrix}.
 \label{eq:r8_increment}
\end{equation}
If the line segment from the current state to the proposed state first meets the boundary at fraction $\alpha<1$, liquidation takes place at time $\alpha h$, at that boundary point. Otherwise $\alpha=1$ and the arrival is settled on the four neighboring grid nodes by independent, mean-preserving two-point lotteries in the two coordinates. Flow utility and the operating benefit accrue for $\alpha h$, with the discount annuity $(1-e^{-\rho\alpha h})/\rho$; liquidation pays $e^{-\rho\alpha h}G$. The complete conventions, including arrival on a boundary node, are given in Supplement S.12.

This is a discrete rebalancing and randomized-settlement contract. The grid is part of its transition technology: it describes admissible recorded balances and preference tiers, not an assertion that interpolation has no economic effect. It is an analytical economic specification, not a calibration of an observed intermediary or an empirical claim that wealth is randomly rounded. The continuous diffusion in Section \ref{sec:preferences} remains a separate economic benchmark. Both specifications preserve consumption, costly preference formation, risky investment, stochastic preferences, covariance, stopping, and liquidation. Their relationship is an approximation question rather than an identity.

An interior arrival $Y$ settled to $Z$ satisfies
\begin{equation}
 \E[Z\mid Y]=Y,\qquad
 \operatorname{Cov}(Z)-\operatorname{Cov}(Y)
 =\diag\left(\E[(Y_i-l_i(Y))(r_i(Y)-Y_i)]\right).
 \label{eq:r8_rounding_variance}
\end{equation}
The quantities $l_i,r_i$ are adjacent grid nodes. Consequently, changing settlement granularity alters risk even when the rebalancing calendar and action menu are held fixed. The wealth-lottery perturbations below exploit this observation: they change an economically relevant primitive and reveal the movement of the original decision region. They are not used to assert an unproved diffusion limit.

\subsection{A Positive-Duration Risk Choice}
Let $\mathcal P^\delta_+$ require the initial risky position to be positive throughout the first commitment interval of length $\delta$, until liquidation, and let $\mathcal P^\delta_-$ require it to be nonpositive. In the finite experiment the control is constant within the first interval, so this is precisely the first-action partition used by the solver. Later positions are unrestricted. In continuous time a constant initial position is one possible commitment contract; an adapted sign restriction over the same interval is a different admissible commitment contract. Either can be studied, but they must not be interchanged without changing the policy sets.

\begin{proposition}[The duration of a risk classification]
\label{prop:r8_duration}
For the diffusion in Section \ref{sec:preferences}, changing a bounded measurable control at a single deterministic initial instant leaves its state path and payoff unchanged, up to the usual indistinguishability convention. Hence a partition that restricts only the sign of that isolated control value has equal class suprema. A strictly signed lifetime class-value difference instead concerns a positive-duration restriction such as $\mathcal P^\delta_\sigma$. When a time mesh is refined while the latter estimand is held fixed, the restriction must continue to cover the same interval $\delta$.
\end{proposition}
\begin{proof}
Changing a bounded drift integrand on a singleton changes no Lebesgue integral. The difference of the corresponding stochastic integrals has quadratic variation equal to the integral of a bounded squared integrand over that singleton, which is zero. The resulting state and stopped payoff therefore coincide. Either initial sign can be assigned without changing the value. Holding the duration of a restriction fixed under refinement is then a condition on the experiment, not an attainment assumption.
\end{proof}

The economic horizon of the reported choice is thus one eighth of a year. This is distinct from the lifetime duration feature $A_p$ in \eqref{eq:risk_class_values}. A local Hamiltonian comparison has flow-value units and can also be informative, but it is not the lifetime commitment-value difference reported here.

\subsection{Preferences over Preference Changes}
\label{subsec:r8_cardinal}
Consumption risk aversion does not identify the cardinal payoff from changing the preference index. A useful explicit specification is
\begin{equation}
 U_b(c,u)=\int_1^c x^{-u}\,dx+b(u)
 =\frac{c^{1-u}-1}{1-u}+b(u).
 \label{eq:r8_metapreference}
\end{equation}
All members have marginal consumption utility $c^{-u}$ and relative risk aversion $u$. The reference consumption level is one in the fixed units of the economy; $b(u)$ is the payoff from being in preference state $u$ at that reference consumption. It is a preference over preference states. It must be specified or, in an empirical application, separately disciplined. The baseline corresponds to $b_0(u)=1/(1-u)$. Thus its cross-preference cardinal comparison is a substantive primitive rather than a consequence of the CRRA coefficient alone.

Consider first $b(u)=b_0(u)+a(u-u_0)$ in the two-stage family of Section \ref{sec:r7_primitive}. A common tilt $a$ changes neither class's consumption lottery nor its fixed-$u_0$ payoff. Nevertheless, it changes the shadow benefit of adjustment in both classes. The following uniform result characterizes a nontrivial set on which the original economic mechanism survives.

\begin{proposition}[A cardinally robust primitive adjustment premium]
\label{prop:r8_cardinal}
Retain $u_0=2$, the adjustment bound $0.2$, $C_-=0.5$, $C_+\in\{0.05,0.8\}$ with probability $p\in[0.06,0.10]$ on the lower outcome, and $k\in[20,80]$. For every common tilt $a\in[-0.25,0.25]$,
\begin{equation}
 \Omega_+(a)-\Omega_-(a)
 \geq \frac{(1.482-0.25)^2\,20/(20+17.342)
                    -(0.6138+0.25)^2}{160}
 >0.0004173.
 \label{eq:r8_cardinal_bound}
\end{equation}
For a duration difference $D_A=0.5$, the adjusted indifference contract is therefore at least $0.0008346$ below the unadjusted one throughout this probability--cost--cardinal family.
\end{proposition}

The proof in Supplement S.12 uses the same curvature restriction as the original proposition and treats the common cost in both option bounds jointly. It is an analytic continuum statement, not an inference from a table of optimized examples. Its sufficient tilt radius is conservative. Conversely, at $(p,k)=(0.08,40)$ the relative option is $0.04882837$ at $a=0$ but $-0.01391977$ at $a=1$. The unadjusted indifference is $1.5$ in both cases; the adjusted indifference changes from $1.40234327$ to $1.52783953$. The larger tilt is outside the robust family, and the reversal is retained as economically informative. Downside exposure and the cardinal value of preference changes jointly determine the mechanism.

\subsection{A Dynamic Accounting and Reoptimization Experiment}
\label{subsec:r8_dynamic}
The dynamic economy does not inherit the two-stage direction mechanically. Even without deliberate adjustment, its preference state is stochastic. We therefore recompute both risk classes in both adjustment regimes under each change in primitives. In particular, an additive tilt need not leave the dynamic no-adjustment values unchanged.

For a policy $p$, let $F_p$ be discounted exposure to $(c^{1-u}-1)/(1-u)$; let $B_p$ and $H_p$ be the corresponding exposures to $1/(1-u)$ and $u-2$; and retain discounted operating duration $A_p$, quadratic adjustment effort $C_p$, and liquidation payoff $G_p$. Under a normalization multiplier $s$, tilt $a$, and liquidation multiplier $z$, the exact policy value is
\begin{equation}
 J_p(s,a,d,k,z)=F_p+sB_p+aH_p+dA_p-kC_p+zG_p.
 \label{eq:r8_feature_account}
\end{equation}
The baseline is $(s,a,k,z)=(1,0,2,1)$. All six features use the stopped transition protocol and include its endogenous exit distribution.

\begin{proposition}[Dynamic feature accounting with reoptimization]
\label{prop:r8_mechanisms}
Fix the finite transition law and the class and adjustment restrictions. The optimized value is the maximum of the affine forms in \eqref{eq:r8_feature_account}. At a differentiability point its derivatives in $(s,a,d,k,z)$ are $(B,H,A,-C,G)$ for an optimizing policy. If $p_0$ is optimal at $q_0=(s_0,a_0,d_0,k_0,z_0)$, then for any $q_1$,
\begin{equation}
 W(q_1)-W(q_0)
 =J_{p_0}(q_1)-J_{p_0}(q_0)+\mathcal R(q_1;p_0),
 \qquad \mathcal R(q_1;p_0)\geq0.
 \label{eq:r8_reoptimization}
\end{equation}
The identity holds separately for all four class--regime values. For the relative option, subtract the two nonpositive-class and no-adjustment comparisons with their appropriate signs; the resulting difference of reoptimization gains has no predetermined sign.
\end{proposition}
\begin{proof}
There are finitely many deterministic Markov policies. The affine representation gives the envelope statement. Add and subtract $J_{p_0}(q_1)$, use $W(q_0)=J_{p_0}(q_0)$, and use feasibility of $p_0$ at $q_1$ to obtain \eqref{eq:r8_reoptimization}. Taking the four signed differences preserves the identity, but not the sign of its remainder.
\end{proof}

Equation \eqref{eq:r8_reoptimization} separates the changed valuation of a fixed exposure from endogenous policy adjustment. It does not identify a causal contribution from data or decompose interacting primitives into invariant shares. For changes in the shock law we directly reevaluate the baseline policy under the new law and again report the nonnegative gain from reoptimization. That calculation changes the feature distribution itself, rather than treating covariance as a reward coefficient. We also solve a zero-covariance law: setting $\lambda=1/2$ would mix two nonzero-correlation shock laws and is not the same intervention.

The results reported below use the original $k=2$ and all financial and preference controls. They explain why a positive relative adjustment option can coexist with positive risk in both regimes, with nonpositive risk in both regimes, or with opposite choices. That separation is the economic content needed to connect a computed certificate to the mechanism of the full model.

\subsection{Decision Error across Economic Targets}
\label{subsec:r8_target_error}
Let $[\ell_r,u_r]$ be a certified difference interval for a finite protocol in regime $r\in\{\mathrm{adj},0\}$, including its numerical arithmetic allowance. If another target satisfies $|W_\sigma^{r,\mathrm{target}}-W_\sigma^{r,\mathrm{finite}}|\leq\eta_\sigma^r$, then
\begin{equation}
 \Delta^{r,\mathrm{target}}\in
 [\ell_r-\eta_+^r-\eta_-^r,\;u_r+\eta_+^r+\eta_-^r].
 \label{eq:r8_target_budget}
\end{equation}
This follows by subtracting the two class-value error intervals. It must not charge an already included arithmetic error a second time. For the historical rectangle, an additional common per-class error strictly below $2.6496827802505486\times10^{-5}$ would suffice for both original signs. That number is a required budget, not an estimate of the diffusion error.

Our robustness experiment instead uses explicitly specified finite settlement protocols and reoptimizes each. It retains all four corners and the center of the historical rectangle, reports its failures as well as its successes, brackets the movement of indifference contracts, and additionally certifies a common nested region. No unexplained constant is substituted into the approximation theorem to turn a finite-model statement into a diffusion claim. The approximation results and continuous-time economy are preserved; the numerical evidence identifies exactly the decisions for which a further target-error bound is consequential.
